% !TEX root = main.tex
%%%%%%%%%%%%%%%%%%%%%%%%%%%%%%%%%%%%%%%%%
% Główny plik pracy
% Szablon pracy dyplomowej
% Wydział Informatyki 
% Zachodniopomorski Uniwersytet Technologiczny w Szczecinie
% autor Joanna Kołodziejczyk (jkolodziejczyk@zut.edu.pl)
% Bardzo wczesnym pierwowzorem szablonu był
% The Legrand Orange Book
% Version 5.0 (29/05/2025)
%
% Modifications to LOB assigned by %JK
%%%%%%%%%%%%%%%%%%%%%%%%%%%%%%%%%%%%%%%%%

%----------------------------------------------------------------------------------------
%	PAKIETY ORAZ PLIKI ZAWIERAJĄCE DEFINICJE STYLI I KONFIGURACJĘ OTOCZEŃ LATEX
%----------------------------------------------------------------------------------------

\documentclass[12pt,fleqn,twoside,openany]{book} % JK Rozmiary czcionek są zgodne z wymaganiami uczelni 
% UWAGA! - wydruk jest dwustronny i jest to zabieg zamierzony

% \usepackage{fix-cm} % Zapewnienie, że wszystkie rozmiary czcionek są dostępne

\usepackage{float}
\input{structure.tex} % JK - Plik zawierający podstawowe elementy konfigurujące układ dokumentu
% UWAGA! - raczej nie będzie potrzeby zmieniania jego struktury

%%%%%%%%%%%%%%%%%%%%%%%%%%%%%%%%%%%%%%%%%
% Plik z definicjami
% Szablon pracy dyplomowej
% Wydział Informatyki 
% Zachodniopomorski Uniwersytet Technologiczny w Szczecinie
% autor Joanna Kołodziejczyk (jkolodziejczyk@zut.edu.pl)
% Bardzo wczesnym pierwowzorem szablonu był
% The Legrand Orange Book
% Version 5.0 (29/05/2025)
%
% Modifications to LOB assigned by %JK
%%%%%%%%%%%%%%%%%%%%%%%%%%%%%%%%%%%%%%%%%



\def\HRule{\color{blueWI} \rule{\linewidth}{0.6pt}} % horisontal rule in ZUT color

%----------------------------------------------------------------------------------------
% Typ pracy (wybrać właściwy)
%----------------------------------------------------------------------------------------
%\def\degreename{Praca dyplomowa magisterska} 
\def\degreename{Praca dyplomowa inżynierska}
%\def\degreename{Praca dyplomowa licencjacka}

%----------------------------------------------------------------------------------------
% Temat pracy
%----------------------------------------------------------------------------------------
\def\ttitle{System wspomagania zarządzania gospodarstwem domowym}
\def\shorttitle{System zarządzania gospodarstwem domowym} %Jeżeli tytuł pracy jest na tyle długi, że zajmuje 3 linijki to trzeba podać krótszy ekwiwalent do nagłówków stron parzystych
\def\ttitleEng{Household Management Support System} %temat pracy w j. angielskim


%----------------------------------------------------------------------------------------
% Informacje o autorze
%----------------------------------------------------------------------------------------
\def\authornames{Imię i Nazwisko} %imię i nazwisko autora
\def\noalbum{XXXXXX}
\def\speciality{Inżynieria Oprogramowania} %nazwa specjalności
\def\field{Informatyka} %dziedzina nauki

%----------------------------------------------------------------------------------------
% Informacje o promotorze
%----------------------------------------------------------------------------------------
\def\supname{dr inż. Imię Nazwisko} %imię i nazwisko promotora
\def\departmentname{Katedra Inżynierii Oprogramowania} %nazwa katedry promotora

%----------------------------------------------------------------------------------------
% Informacje o promotorze dodatkowym (odkomentować w razie potrzeby)
%----------------------------------------------------------------------------------------
% \def\supExt{prof. dr hab. inż. Marian Kowalski} %imię i nazwisko promotora
% \def\departmentnameExt{Katedra Grafiki Komputerowej} %nazwa katedry promotora

%----------------------------------------------------------------------------------------
% Data wydania tematu pracy
%----------------------------------------------------------------------------------------
\def\datetitle{01.10.2025}

%----------------------------------------------------------------------------------------
% Rok i miejsce złożenia pracy
%----------------------------------------------------------------------------------------
\def\placesubmit{Szczecin}
\def\yearsubmit{2026}
  % JK - dodatkowe definicje głównie treść strony tytułowej
% UWAGA! - konieczność edycji celem zmiany autora/tematu/dat itp., itd


%----------------------------------------------------------------------------------------
% OTWARCIE DOKUMENTU
%----------------------------------------------------------------------------------------
\begin{document}

%----------------------------------------------------------------------------------------
% STRONA TYTUŁOWA 
%----------------------------------------------------------------------------------------
\include{title_page}

%----------------------------------------------------------------------------------------
% PUSTA STRONA PO STRONIE TYUŁOWEJ (JK - design and implementation)
%----------------------------------------------------------------------------------------
\newpage
\thispagestyle{empty}
~\vfill

%----------------------------------------------------------------------------------------
%	STRESZCZENIE I SŁOWA KLUCZOWE (1 STRONA) (JK - design and implementation)
%----------------------------------------------------------------------------------------
\newpage
\thispagestyle{empty}
%%%%%%%%%%%%%%%%%%%%%%%%%%%%%%%%%%%%%%%%%
% Specjalna strona pracy ze streszczeniem i abstractem w j. angielskim
% Szablon pracy dyplomowej
% Wydział Informatyki 
% Zachodniopomorski Uniwersytet Technologiczny w Szczecinie
% autor Joanna Kołodziejczyk (jkolodziejczyk@zut.edu.pl)
% Bardzo wczesnym pierwowzorem szablonu był
% The Legrand Orange Book
% Version 5.0 (29/05/2025)
%
% Modifications to LOB assigned by %JK
%%%%%%%%%%%%%%%%%%%%%%%%%%%%%%%%%%%%%%%%%


\begin{center}
\noindent {{\color{blueZUT}\Large\sffamily  {Streszczenie}}}\\[1cm] 
\end{center}

   Przedmiotem niniejszej pracy jest projekt i implementacja systemu informatycznego ,,Household Manager'' wspomagającego zarządzanie gospodarstwem domowym oraz zintegrowanymi urządzeniami inteligentnego domu (smart home). Aplikacja rozwiązuje problemy związane z podziałem obowiązków, rozliczaniem wspólnych wydatków oraz koordynacją działań między domownikami.

System składa się z aplikacji mobilnej stworzonej w technologii React Native oraz serwera backendowego napisanego w języku Go, wykorzystującego bazę danych PostgreSQL. Kluczowe funkcjonalności obejmują: zarządzanie wirtualnymi domami, moduł sterowania urządzeniami smart home, system przydzielania zadań, moduł finansowy z automatycznym odczytywaniem danych z rachunków (OCR), wspólne listy zakupów oraz ankiety decyzyjne. Zastosowanie nowoczesnych technologii pozwoliło na stworzenie wydajnego, skalowalnego i intuicyjnego w obsłudze rozwiązania, dostępnego na platformy Android i iOS.

\vspace{10pt}
\noindent{\bf słowa kluczowe:} zarządzanie gospodarstwem domowym, inteligentny dom, aplikacja mobilna, podział obowiązków, wspólne wydatki, React Native, Go

\vfill

\begin{center}
\noindent {{\color{blueZUT}\Large\sffamily {Abstract}}}\\[1cm] 
\end{center}

The subject of this thesis is the design and implementation of the “Household Manager” IT system, which supports household management and integrated smart home devices. The application solves problems related to the division of responsibilities, settlement of joint expenses, and coordination of activities between household members.

The system consists of a mobile application created in React Native technology and a backend server written in Go, using a PostgreSQL database. Key features include: virtual home management, a smart home device control module, a task assignment system, a financial module with automatic bill reading (OCR), shared shopping lists, and decision-making surveys. The use of modern technologies has allowed us to create an efficient, scalable, and intuitive solution, available on Android and iOS platforms.

\vspace{10pt}
\noindent{\bf keywords:} household management, smart home, mobile application, task management, shared expenses, React Native, Go
 %abstract.tex 
 
%----------------------------------------------------------------------------------------
%	SPIS TREŚCI (JK - design and implementation)
%----------------------------------------------------------------------------------------
\pagestyle{empty}  % Wyłącz stopkę i nagłówek w TOC
\tableofcontents % wyświetl spis
\addtocontents{toc}{\protect\thispagestyle{empty}} % Zachowaj pusty nagłówek i stopkę w TOC
% Wymuszenie rozpoczęcie pierwszego rozdziału na nieparzystej stronie, aby znajdował się po prawej stronie 
\cleardoublepage 
% Ponownie włącz nagłówki i stopki
\pagestyle{fancy} 

%----------------------------------------------------------------------------------------
%	WSTĘP W PRACY DYPLOMOWEJ (JK - design and implementation)
%----------------------------------------------------------------------------------------
\addcontentsline{toc}{chapter}{Wstęp}
%%%%%%%%%%%%%%%%%%%%%%%%%%%%%%%%%%%%%%%%%
% Plik z wstępem do pracy
% Szablon pracy dyplomowej
% Wydział Informatyki 
% Zachodniopomorski Uniwersytet Technologiczny w Szczecinie
% autor Joanna Kołodziejczyk (jkolodziejczyk@zut.edu.pl)
% Bardzo wczesnym pierwowzorem szablonu był
% The Legrand Orange Book
% Version 5.0 (29/05/2025)
%
% Modifications to LOB assigned by %JK
%%%%%%%%%%%%%%%%%%%%%%%%%%%%%%%%%%%%%%%%%


\chapter*{Wstęp}

Współczesne tempo życia oraz rosnąca liczba obowiązków sprawiają, że efektywne zarządzanie gospodarstwem domowym staje się wyzwaniem, szczególnie w przypadku osób dzielących przestrzeń życiową, takich jak studenci, współlokatorzy czy rodziny. Codzienne zadania, takie jak sprzątanie, zakupy, płacenie rachunków czy podejmowanie wspólnych decyzji, często stają się źródłem nieporozumień i konfliktów, gdy brakuje jasnego podziału odpowiedzialności i sprawnej komunikacji.

Tradycyjne metody organizacji, takie jak kartki na lodówce, ustne ustalenia czy chaotyczne konwersacje w komunikatorach internetowych, okazują się niewystarczające w obliczu złożoności współczesnego życia. Brak centralnego systemu do śledzenia wydatków, terminów płatności czy postępów w realizacji zadań domowych prowadzi do zapominania o ważnych terminach, dublowania zakupów oraz poczucia niesprawiedliwości w podziale obowiązków.

Celem niniejszej pracy jest zaprojektowanie i zaimplementowanie kompleksowego systemu informatycznego wspomagającego zarządzanie gospodarstwem domowym. Aplikacja ma na celu ułatwienie komunikacji między domownikami, automatyzację procesów związanych z płatnościami i zakupami oraz zapewnienie przejrzystości w podziale obowiązków.

Głównym elementem systemu jest aplikacja mobilna, która umożliwia użytkownikom tworzenie wirtualnych domów, zapraszanie do nich współlokatorów oraz wspólne zarządzanie zasobami. Kluczowe funkcjonalności obejmują system zarządzania zadaniami z możliwością przypisywania ich do konkretnych osób, moduł finansowy do śledzenia rachunków i wydatków (wspierany technologią OCR do automatycznego odczytu danych z faktur), wspólne listy zakupów oraz system ankiet do podejmowania demokratycznych decyzji.

Praca składa się z sześciu rozdziałów. Rozdział pierwszy zawiera analizę problemu oraz przegląd istniejących rozwiązań. Rozdział drugi przedstawia specyfikację wymagań funkcjonalnych i niefunkcjonalnych oraz projekt systemu. Rozdział trzeci opisuje wybrane technologie wykorzystane do realizacji projektu, w tym język Go dla warstwy serwerowej oraz framework React Native dla aplikacji mobilnej. Rozdział czwarty poświęcony jest szczegółom implementacyjnym, w tym strukturze bazy danych, architekturze API oraz rozwiązaniom z zakresu bezpieczeństwa. Rozdział piąty prezentuje wyniki testów i badań nad wytworzonym oprogramowaniem. Całość kończy podsumowanie i wnioski z realizacji projektu. % introduction.tex zawiera treść wstępu

%----------------------------------------------------------------------------------------
%	ROZDZIAŁ 1
%----------------------------------------------------------------------------------------
\chapter{Analiza problemu i cel projektu}
\label{rozdzial1}

Wspólne mieszkanie, czy to w gronie rodziny, czy ze współlokatorami, wiąże się z koniecznością koordynacji wielu aspektów życia codziennego. Rozdział ten poświęcony jest analizie problemów związanych z zarządzaniem gospodarstwem domowym, zdefiniowaniu celu projektowanej aplikacji oraz przeglądowi istniejących na rynku rozwiązań.

\section{Opis problemu współdzielenia obowiązków domowych}

Zarządzanie domem to złożony proces, na który składa się szereg powtarzalnych czynności oraz decyzji podejmowanych przez grupę osób. Do najczęstszych problemów, z jakimi borykają się domownicy, należą:

\begin{itemize}
    \item \textbf{Niejasny podział obowiązków:} Brak precyzyjnego przypisania zadań (np. sprzątanie łazienki, wynoszenie śmieci) prowadzi do sytuacji, w których obowiązki są zaniedbywane lub wykonywane ciągle przez te same osoby, co rodzi frustrację i konflikty.
    \item \textbf{Problemy z rozliczeniami finansowymi:} Wspólne wydatki na czynsz, media czy zakupy spożywcze wymagają sprawiedliwego podziału kosztów. Ręczne spisywanie wydatków jest uciążliwe i podatne na błędy, a brak terminowych spłat między współlokatorami zaburza płynność finansową.
    \item \textbf{Dublowanie zakupów lub braki w zapasach:} Brak zsynchronizowanej listy zakupów skutkuje kupowaniem niepotrzebnych produktów lub - przeciwnie - zapominaniem o artykułach pierwszej potrzeby.
    \item \textbf{Trudności w komunikacji i podejmowaniu decyzji:} Ustalenia dotyczące wspólnych spraw (np. wybór dostawcy internetu, planowanie posiłków) często giną w natłoku wiadomości na komunikatorach, co utrudnia osiągnięcie konsensusu.
\end{itemize}

Brak zorganizowanego systemu wsparcia tych procesów obciąża domowników tzw. ,,obciążeniem mentalnym'', polegającym na konieczności ciągłego pamiętania o terminach i zobowiązaniach.

\section{Cel i zakres aplikacji}

Głównym celem projektu jest stworzenie aplikacji mobilnej ,,Household Manager'', która zintegruje kluczowe funkcjonalności niezbędne do sprawnego zarządzania domem w jednym, intuicyjnym interfejsie. System ma za zadanie zminimalizować ryzyko konfliktów poprzez zapewnienie transparentności działań i automatyzację powtarzalnych czynności.

Zakres funkcjonalny aplikacji obejmuje:
\begin{enumerate}
    \item \textbf{Zarządzanie domami i użytkownikami:} Możliwość tworzenia wirtualnego ,,domu'', zapraszania użytkowników poprzez kody lub linki oraz zarządzanie ich uprawnieniami.
    \item \textbf{System zadań (Tasks):} Tworzenie, przypisywanie i śledzenie statusu prac domowych.
    \item \textbf{Moduł finansowy (Bills):} Rejestrowanie rachunków, podział kosztów (równy lub według wag), śledzenie statusu płatności. Istotnym elementem jest integracja z mechanizmem OCR (Optical Character Recognition) pozwalającym na automatyczne wprowadzanie danych z fizycznych rachunków i faktur.
    \item \textbf{Listy zakupów (Shopping Lists):} Współdzielone listy produktów z aktualizacją w czasie rzeczywistym, co pozwala na dodawanie i odhaczanie pozycji przez wszystkich domowników.
    \item \textbf{System ankiet (Polls):} Narzędzie do demokratycznego podejmowania decyzji grupowych.
    \item \textbf{Zarządzanie inteligentnym domem (Smart Home):} Integracja, sterowanie oraz monitorowanie podłączonych urządzeń domowych z poziomu jednej aplikacji, co centralizuje kontrolę nad gospodarstwem domowym.
    \item \textbf{Powiadomienia:} System alertów przypominających o nadchodzących terminach płatności czy przydzielonych zadaniach.
\end{enumerate}

Aplikacja jest kierowana głównie do osób młodych, studentów wynajmujących wspólnie mieszkania oraz nowoczesnych rodzin, dla których smartfon jest naturalnym narzędziem organizacji życia.

\section{Analiza istniejących rozwiązań}

Na rynku dostępnych jest wiele narzędzi wspierających organizację, jednak większość z nich skupia się na wąskich dziedzinach lub jest adaptacją narzędzi biznesowych.

\subsection{Splitwise}
Splitwise to popularna aplikacja do dzielenia wydatków. Doskonale radzi sobie z matematyką podziału kosztów i upraszczaniem długów.
\begin{itemize}
    \item \textbf{Zalety:} Rozbudowane opcje podziału kosztów, obsługa wielu walut.
    \item \textbf{Wady:} Brak funkcji zarządzania zadaniami domowymi, listami zakupów czy komunikacji grupowej niezwiązanej z finansami.
\end{itemize}

\subsection{Trello / Asana}
Narzędzia typu Kanban, popularne w zarządzaniu projektami w IT.
\begin{itemize}
    \item \textbf{Zalety:} Elastyczność, możliwość tworzenia dowolnych tablic i list zadań.
    \item \textbf{Wady:} Zbyt skomplikowane do prostych zastosowań domowych, brak dedykowanych funkcji finansowych (np. podział rachunków), interfejs nastawiony na workflow biznesowy.
\end{itemize}

\subsection{OurGroceries}
Aplikacja dedykowana listom zakupów.
\begin{itemize}
    \item \textbf{Zalety:} Prosta obsługa, synchronizacja w czasie rzeczywistym.
    \item \textbf{Wady:} Ograniczona tylko do zakupów, przestarzały interfejs użytkownika (UI), brak integracji z finansami i zadaniami.
\end{itemize}

\subsection{Rozwiązania zintegrowane (np. OurHome, Cozi)}
Istnieją aplikacje próbujące łączyć te funkcje, jednak często borykają się z problemami:
\begin{itemize}
    \item Część z nich jest płatna w modelu subskrypcyjnym, co zniechęca studentów.
    \item Wiele projektów nie jest już rozwijanych i posiada przestarzały design.
    \item Rzadko oferują zaawansowane funkcje jak skanowanie rachunków (OCR).
\end{itemize}

\section{Podsumowanie analizy}

Analiza rynku wskazuje na brak darmowego, nowoczesnego i kompleksowego rozwiązania, które łączyłoby zarządzanie finansami (z wykorzystaniem OCR), zadaniami i zakupami w jednej aplikacji o wysokiej jakości interfejsu (UX/UI). Projektowany system ,,Household Manager'' ma na celu wypełnienie tej luki, oferując użytkownikom narzędzie typu ,,wszystko w jednym'', dostosowane do współczesnych standardów aplikacji mobilnych.

\section{Metodyka pracy i zarządzanie projektem}

W celu zapewnienia wysokiej efektywności prac oraz terminowej realizacji wszystkich założeń projektowych, proces wytwarzania oprogramowania był prowadzony zgodnie z metodyką zwinnego zarządzania projektami (ang. \textit{Agile}).

\subsection{Zastosowanie podejścia Kanban}

Do zarządzania zadaniami wybrano podejście \textbf{Kanban}, które charakteryzuje się wizualizacją przepływu pracy oraz skupieniem na ciągłym dostarczaniu kolejnych funkcjonalności. Wybór ten był podyktowany elastycznością metody, która pozwala na dynamiczne reagowanie na napotkane problemy techniczne oraz płynne przechodzenie między etapami implementacji warstwy serwerowej i klienckiej.

\subsection{Narzędzia wspomagające --- GitHub Projects}

Głównym narzędziem wykorzystywanym do zarządzania cyklem życia projektu była platforma \textbf{GitHub}, a w szczególności usługa \textbf{GitHub Projects}. Wszystkie wymagania funkcjonalne zostały zdekomponowane na mniejsze, atomowe jednostki pracy (ang. \textit{Issues}), które następnie trafiły do rejestru zadań produktu (ang. \textit{Backlog}).

Proces pracy został odzwierciedlony na cyfrowej tablicy Kanban składającej się z kolumn:
\begin{itemize}
    \item \textbf{Todo (Do zrobienia):} Zadania zaplanowane do realizacji, czekające w kolejce.
    \item \textbf{In Progress (W toku):} Funkcjonalności będące w danej chwili implementowane.
    \item \textbf{In Review (W weryfikacji):} Zadania ukończone programistycznie, przechodzące etap testów manualnych lub automatycznych.
    \item \textbf{Done (Ukończone):} W pełni przetestowane i wdrożone moduły aplikacji.
\end{itemize}

\begin{figure}[H]
    \centering
    \includegraphics[width=0.9\textwidth]{backlog.png}
    \caption{Tablica Kanban projektu w systemie GitHub Projects}
    \label{fig:backlog}
\end{figure}

Takie podejście umożliwiło precyzyjne śledzenie postępów prac oraz zapewniło transparentność całego procesu, co jest kluczowe przy realizacji złożonych systemów typu full-stack. Dzięki powiązaniu zadań z systemem kontroli wersji Git, każda zmiana w kodzie była bezpośrednio mapowana na konkretne wymaganie biznesowe zawarte w projekcie. % chapter1.tex zawiera treść rozdziału 1

%----------------------------------------------------------------------------------------
%	ROZDZIAŁ 2
%----------------------------------------------------------------------------------------
\chapter{Specyfikacja wymagań i projekt systemu}
\label{rozdzial2}

Rozdział ten definiuje wymagania stawiane systemowi ,,Household Manager'' oraz przedstawia jego projekt architektoniczny i bazodanowy.

\section{Wymagania funkcjonalne}

Wymagania funkcjonalne określają, co system powinien robić. Zostały one podzielone na moduły odpowiadające głównym obszarom działania aplikacji.

\subsection{Moduł Użytkownika i Uwierzytelniania}
\begin{itemize}
    \item \textbf{Rejestracja i logowanie:} Użytkownik może założyć konto podając adres e-mail i hasło lub skorzystać z logowania przez Google (OAuth 2.0). System weryfikuje poprawność wprowadzonych danych (np. format adresu e-mail, złożoność hasła).
    \item \textbf{Zarządzanie profilem:} Możliwość edycji awatara (zdjęcia profilowego) oraz danych osobowych (imię, nazwisko). Zmiana hasła wymaga ponownego uwierzytelnienia.
    \item \textbf{Przypomnienie hasła:} Procedura resetowania hasła drogą mailową, wykorzystująca unikalne tokeny jednorazowe.
\end{itemize}

\subsection{Moduł Domu (Home)}
\begin{itemize}
    \item \textbf{Tworzenie domu:} Użytkownik może utworzyć nową przestrzeń (Dom), nadając jej nazwę. Automatycznie staje się administratorem nowo utworzonego domu.
    \item \textbf{Zarządzanie członkami:} Możliwość zapraszania nowych osób poprzez unikalny, 6-cyfrowy kod zaproszenia, ważny przez określony czas. Administrator może usuwać członków z domu.
    \item \textbf{Wiele domów:} Architektura systemu pozwala użytkownikowi na przynależność do wielu domów jednocześnie (np. "Mieszkanie Studenckie" i "Dom Rodzinny") oraz płynne przełączanie się między nimi w interfejsie aplikacji.
\end{itemize}

\subsection{Moduł Zadań (Tasks)}
\begin{itemize}
    \item \textbf{CRUD Zadań:} Pełna obsługa cyklu życia zadania: tworzenie, edycja, usuwanie i wyświetlanie listy zadań. Każde zadanie posiada tytuł, opis oraz priorytet.
    \item \textbf{Przydzielanie:} Możliwość przypisania zadania do jednego lub wielu domowników. System obsługuje również zadania rotacyjne (cykliczne).
    \item \textbf{Terminy:} Definiowanie daty wykonania zadania (deadline). Aplikacja wyróżnia zadania przeterminowane.
    \item \textbf{Statusy:} Oznaczanie zadań jako wykonane, co jest widoczne dla wszystkich członków domu w czasie rzeczywistym.
\end{itemize}

\subsection{Moduł Finansowy (Bills)}
\begin{itemize}
    \item \textbf{Dodawanie rachunków:} Wprowadzanie kwoty, tytułu, kategorii (np. media, czynsz, zakupy) oraz terminu płatności.
    \item \textbf{OCR (Optical Character Recognition):} Innowacyjna funkcja pozwalająca na zrobienie zdjęcia fizycznego paragonu lub faktury. System automatycznie przetwarza obraz, ekstrahując kwotę całkowitą i nazwę sprzedawcy, co znacząco przyspiesza proces wprowadzania danych.
    \item \textbf{Podział kosztów:} Elastyczny podział kwoty rachunku na członków domu. Domyślnie system proponuje podział równomierny, ale użytkownik może go modyfikować.
    \item \textbf{Rozliczanie:} Śledzenie statusu płatności (kto już zapłacił swoją część, a kto zalega).
\end{itemize}

\subsection{Moduł Zakupów (Shopping)}
\begin{itemize}
    \item \textbf{Lista zakupów:} Dodawanie produktów do wspólnej listy zakupowej.
    \item \textbf{Kategorie:} Grupowanie produktów (np. spożywcze, chemia, warzywa) w celu łatwiejszej nawigacji w sklepie.
    \item \textbf{Interakcja:} Odhaczanie kupionych produktów, co skutkuje ich przeniesieniem do sekcji "Kupione" lub usunięciem z listy. Synchronizacja stanu listy między urządzeniami zapobiega dublowaniu zakupów.
\end{itemize}

\subsection{Moduł Ankiet (Polls)}
\begin{itemize}
    \item \textbf{Tworzenie ankiet:} Możliwość zadawania pytań dotyczących wspólnych spraw (np. "Co na obiad?", "Jakiego dostawcę internetu wybieramy?") i definiowania opcji odpowiedzi.
    \item \textbf{Głosowanie:} Anonimowe lub jawne oddawanie głosu przez domowników.
    \item \textbf{Wyniki:} Wizualizacja wyników ankiety w formie wykresu lub listy procentowej.
\end{itemize}

\section{Wymagania niefunkcjonalne}

\begin{itemize}
    \item \textbf{Dostępność (Availability):} System powinien być dostępny 24 godziny na dobę, 7 dni w tygodniu, z planowanymi oknami serwisowymi w godzinach nocnych.
    \item \textbf{Skalowalność:} Architektura backendu musi umożliwiać łatwe skalowanie poziome (dodawanie kolejnych instancji serwera) w przypadku wzrostu liczby użytkowników, co realizowane jest poprzez konteneryzację (Docker).
    \item \textbf{Bezpieczeństwo:}
        \begin{itemize}
            \item Hasła użytkowników muszą być przechowywane w formie zahaszowanej przy użyciu silnego algorytmu (np. bcrypt).
            \item Komunikacja między klientem a serwerem musi być szyfrowana protokołem TLS/SSL (HTTPS).
            \item Dostęp do zasobów API jest zabezpieczony mechanizmem tokenów JWT (JSON Web Tokens).
        \end{itemize}
    \item \textbf{Responsywność i wydajność:} Aplikacja mobilna powinna działać płynnie na systemach Android (wersja 10+) i iOS (wersja 15+). Czas odpowiedzi API nie powinien przekraczać 200ms dla 95\% standardowych zapytań (np. pobranie listy zadań).
    \item \textbf{Internacjonalizacja:} System powinien być przygotowany do obsługi wielu wersji językowych (i18n), z domyślnym wsparciem dla języka polskiego i angielskiego.
\end{itemize}

\section{Projekt systemu}

\subsection{Architektura klient-serwer}

System ,,Household Manager'' został zaprojektowany w oparciu o klasyczną, trójwarstwową architekturę klient-serwer (ang. \textit{Three-Tier Architecture}), która zapewnia logiczną separację warstwy prezentacji, logiki biznesowej oraz danych. Takie podejście ułatwia niezależny rozwój poszczególnych komponentów oraz ich utrzymanie.

\begin{itemize}
    \item \textbf{Warstwa Prezentacji (Klient):} Aplikacja mobilna zaimplementowana w frameworku React Native. Odpowiada za interakcję z użytkownikiem, wyświetlanie danych oraz przesyłanie żądań do serwera. Klient jest "cienki" (ang. \textit{thin client}), co oznacza, że większość ciężkiej logiki biznesowej została przeniesiona na serwer.
    \item \textbf{Warstwa Logiki Biznesowej (Serwer):} Serwer REST API napisany w języku Go. Przetwarza żądania przychodzące z aplikacji mobilnej, wykonuje operacje na danych, autoryzuje użytkowników i komunikuje się z bazą danych.
    \item \textbf{Warstwa Danych (Baza Danych):} Relacyjna baza danych PostgreSQL przechowująca wszystkie trwałe informacje systemu.
\end{itemize}

Poniższy schemat (Rysunek \ref{fig:architecture}) prezentuje ogólną architekturę systemu oraz przepływ danych między komponentami.

\begin{figure}[h]
    \centering
    
\includegraphics[width=0.9\textwidth]{placeholder.jpg} % Placeholder
    \caption{Schemat architektury systemu Household Manager}
    \label{fig:architecture}
\end{figure}

\subsection{Projekt bazy danych PostgreSQL}

Jako system zarządzania bazą danych (RDBMS) wybrano \textbf{PostgreSQL}. Decyzja ta podyktowana była kilkoma kluczowymi czynnikami:

\begin{enumerate}
    \item \textbf{Integralność danych (ACID):} PostgreSQL jest w pełni zgodny ze standardem ACID (Atomicity, Consistency, Isolation, Durability), co jest kluczowe w systemach finansowych (moduł rachunków), gdzie nie można pozwolić sobie na utratę lub niespójność danych transakcyjnych.
    \item \textbf{Wsparcie dla JSONB:} Baza ta oferuje doskonałą obsługę typu danych JSON, co pozwala na elastyczne przechowywanie mniej ustrukturyzowanych danych (np. metadane logów OCR) obok tradycyjnych tabel relacyjnych.
    \item \textbf{Wydajność i rozszerzalność:} PostgreSQL słynie z wysokiej wydajności przy złożonych zapytaniach oraz możliwości tworzenia własnych indeksów i funkcji.
\end{enumerate}

\subsection{Diagram ERD}

Struktura bazy danych została zaprojektowana tak, aby odzwierciedlać rzeczywiste relacje panujące w gospodarstwie domowym. Centralnym punktem schematu są tabele użytkowników (\texttt{users}) oraz domów (\texttt{homes}), połączone relacją wiele-do-wielu poprzez tabelę łącznikową (\texttt{home\_memberships}).

Poniższy diagram ERD (Entity Relationship Diagram) przedstawia szczegółowy schemat bazy danych.

\begin{figure}[h]
    \centering
    
\includegraphics[width=1.0\textwidth]{placeholder.jpg} % Placeholder for ERD
    \caption{Diagram ERD bazy danych}
    \label{fig:erd}
\end{figure}

Kluczowe encje w systemie to:

\begin{itemize}
    \item \textbf{users:} Przechowuje unikalne dane identyfikacyjne użytkowników (adres email, hash hasła) oraz dane profilowe.
    \item \textbf{homes:} Reprezentuje pojedyncze gospodarstwo domowe. Posiada unikalny identyfikator oraz kod zaproszenia.
    \item \textbf{home\_memberships:} Tabela asocjacyjna definiująca przynależność użytkownika do domu oraz jego rolę (np. ADMIN, MEMBER).
    \item \textbf{tasks:} Tabela zadań zawierająca tytuł, opis, datę wykonania oraz status.
    \item \textbf{task\_assignments:} Definiuje, który użytkownik jest odpowiedzialny za dane zadanie. Umożliwia przypisanie wielu osób do jednego zadania.
    \item \textbf{bills:} Przechowuje informacje o rachunkach: kwotę całkowitą, datę płatności, wystawcę rachunku oraz status (opłacony/nieopłacony).
    \item \textbf{shopping\_items} oraz \textbf{shopping\_categories:} Tabele obsługujące moduł zakupów, pozwalające na kategoryzację produktów.
\end{itemize}

\subsection{Modele GORM i relacje}

W warstwie aplikacji (backend w Go) wykorzystano bibliotekę \textbf{GORM} (Go Object Relational Mapping). Jest to potężne narzędzie mapowania obiektowo-relacyjnego, które pozwala programiście operować na strukturach języka Go (structs) zamiast pisać surowe zapytania SQL.

Zalety wykorzystania GORM w projekcie:
\begin{itemize}
    \item \textbf{Bezpieczeństwo:} Automatyczna ochrona przed atakami typu SQL Injection dzięki parametryzacji zapytań.
    \item \textbf{Automigracje:} GORM potrafi automatycznie tworzyć i aktualizować schemat bazy danych na podstawie definicji struktur w kodzie (Code-First approach), co znacznie przyspiesza prototypowanie.
    \item \textbf{Abstrakcja:} Pozwala na łatwą zmianę silnika bazy danych w przyszłości (np. na MySQL czy SQLite) bez konieczności przepisywania logiki aplikacji.
\end{itemize}

Wadą rozwiązań typu ORM może być narzut wydajnościowy przy bardzo skomplikowanych zapytaniach analitycznych, jednak w przypadku typowych operacji CRUD (Create, Read, Update, Delete) realizowanych w systemie Household Manager, korzyści wynikające z szybkości wytwarzania oprogramowania przeważają nad minimalnym spadkiem wydajności. % chapter2.tex zawiera treść rozdziału 2

%----------------------------------------------------------------------------------------
%	ROZDZIAŁ 3
%----------------------------------------------------------------------------------------
\chapter{Wybór technologii}
\label{rozdzial3}

Wybór odpowiednich technologii jest kluczowym etapem procesu wytwarzania oprogramowania, mającym bezpośredni wpływ na wydajność, skalowalność, bezpieczeństwo oraz łatwość późniejszego utrzymania systemu (maintainability). Decyzje podjęte na tym etapie determinują szybkość procesu deweloperskiego oraz komfort pracy programistów. W niniejszym rozdziale przedstawiono szczegółową charakterystykę stosu technologicznego wykorzystanego do realizacji projektu ,,Household Manager''.

\section{Backend — Go, GORM, PostgreSQL, Chi Router}

Warstwa serwerowa (backend) stanowi mózg aplikacji, odpowiadający za logikę biznesową, przetwarzanie danych oraz komunikację z bazą danych. Do jej implementacji wybrano język \textbf{Go} (znany również jako Golang).

\subsection{Uzasadnienie wyboru Go}

Wybór języka Go (Golang) jako fundamentu warstwy serwerowej podyktowany był potrzebą zapewnienia wysokiej wydajności oraz niezawodności przy jednoczesnej prostocie implementacji logiki biznesowej. Jako język nowoczesny, rozwijany przez firmę Google, Go adresuje problemy znane ze starszych technologii, takie jak długi czas kompilacji czy skomplikowane zarządzanie pamięcią.


Główne powody wyboru tej technologii to:

\begin{enumerate}
    \item \textbf{Wydajność kompilowana:} Go jest językiem kompilowanym bezpośrednio do kodu maszynowego, co sprawia, że aplikacje napisane w Go są znacznie szybsze niż te uruchamiane w środowiskach interpretowanych (jak Python czy JavaScript/Node.js). Nie wymaga maszyny wirtualnej (JVM), co przekłada się na mniejsze zużycie pamięci RAM i szybszy start aplikacji.
    
    \item \textbf{Statyczne typowanie:} Silny system typów pomaga wykrywać błędy już na etapie kompilacji, zanim kod trafi na produkcję. W przeciwieństwie do języków dynamicznych, programista ma pewność, że zmienna zadeklarowana jako liczba całkowita zawsze nią będzie, co zwiększa stabilność systemu.
    
    \item \textbf{Współbieżność (Concurrency):} Go posiada unikalny model współbieżności oparty na tzw. \textit{goroutines} i kanałach (\textit{channels}). Goroutine to "lekki wątek", zarządzany przez runtime języka, a nie przez system operacyjny. Uruchomienie tysięcy goroutines wymaga ułamka zasobów potrzebnych dla tradycyjnych wątków systemowych. W kontekście serwera HTTP oznacza to możliwość jednoczesnej obsługi ogromnej liczby zapytań od użytkowników przy minimalnym narzucie sprzętowym.
    
    \item \textbf{Prostota i czytelność:} Składnia Go jest minimalistyczna (brak klas, dziedziczenia w tradycyjnym sensie, skomplikowanych meta-programowania). Wymusza to pisanie prostego, czytelnego kodu, który jest łatwy do zrozumienia dla nowych członków zespołu. Sformatowany kod narzędziem \texttt{gofmt} wygląda zawsze tak samo, co eliminuje spory o styl kodowania.
\end{enumerate}

W porównaniu do popularnego Node.js, Go oferuje lepszą wydajność obliczeniową (istotne np. przy przetwarzaniu obrazów w module OCR) oraz stabilniejszy ekosystem bibliotek standardowych.

\subsection{Framework i biblioteki backendowe}

W ekosystemie Go panuje tendencja do unikania rozbudowanych frameworków na rzecz komponowania aplikacji z mniejszych, wyspecjalizowanych bibliotek. W projekcie wykorzystano następujące rozwiązania:

\begin{itemize}
    \item \textbf{Chi Router:}
    Lekki i szybki router HTTP, w pełni kompatybilny ze standardową biblioteką \texttt{net/http}. Wykorzystano go do budowy przejrzystego systemu \textit{middleware}, odpowiadającego za logowanie żądań, obsługę polityki CORS oraz weryfikację tokenów JWT. Chi zapewnia wysoką czytelność kodu przy minimalnym narzucie wydajnościowym.
    
    \item \textbf{GORM (Go Object Relational Mapping):}
    Biblioteka typu ORM, która umożliwiła mapowanie struktur języka Go na relacje w bazie PostgreSQL. GORM został wykorzystany do automatyzacji operacji CRUD, obsługi złożonych relacji między encjami (np. \textit{many-to-many} dla członków domu) oraz zarządzania migracjami schematu bazy danych.
    
    \item \textbf{Google Gemini API (Generative AI):}
    Zamiast klasycznych silników OCR, w systemie wykorzystano zaawansowany model multimodalny \textit{Gemini 1.5 Flash}. Pozwala on nie tylko na odczyt tekstu (OCR), ale również na jego semantyczną analizę i strukturyzację danych w formacie JSON.
    
    
    
    \item \textbf{Google Cloud SDK / HTTP Client:}
    Do komunikacji z modelem Gemini wykorzystano natywny klient HTTP języka Go. Implementacja (usługa \texttt{OCRService}) obejmuje bezpieczne przesyłanie obrazów zakodowanych w standardzie Base64 oraz autorski algorytm \texttt{calculateGeminiConfidence}, który ocenia wiarygodność zwróconych danych (np. poprzez weryfikację sumy pozycji paragonu z kwotą całkowitą).
\end{itemize}
W pracy dyplomowej najlepiej wygląda układ, w którym każda kluczowa technologia ma swój dedykowany podpunkt. Pozwala to recenzentowi szybko odnaleźć konkretne narzędzia i ocenić, czy student rozumie ich rolę w projekcie.

Oto kompletna, rozbudowana sekcja o froncie, przygotowana w sposób najbardziej "dyplomowy" (oddzielne podpunkty dla każdego narzędzia).

\section{Frontend — React Native, Expo, TypeScript}

Warstwa kliencka (aplikacja mobilna) została zrealizowana w architekturze opartej na komponentach, co zapewnia modularność kodu i łatwość rozbudowy interfejsu użytkownika.

\subsection{React Native i architektura komponentowa}

React Native, stworzony przez firmę Meta, jest frameworkiem pozwalającym na budowanie natywnych aplikacji mobilnych przy użyciu biblioteki React. W przeciwieństwie do rozwiązań hybrydowych, React Native nie korzysta z widoków przeglądarkowych (\textit{WebView}), lecz mapuje komponenty JavaScript na natywne widoki systemów iOS i Android. Zapewnia to wysoką wydajność interfejsu oraz dostęp do płynnych animacji systemowych. Wykorzystanie paradygmatu deklaratywnego UI pozwala na automatyczne aktualizowanie widoku w odpowiedzi na zmiany stanu aplikacji.

\subsection{Język TypeScript}

W projekcie zdecydowano się na użycie języka \textbf{TypeScript} jako nadzbioru języka JavaScript. TypeScript wprowadza statyczne typowanie, co w kontekście pracy inżynierskiej przynosi następujące korzyści:
\begin{itemize}
\item \textbf{Wykrywanie błędów w czasie kompilacji:} Eliminuje typowe błędy typu \textit{runtime}, takie jak próba dostępu do pola w obiekcie o wartości \textit{null}.
\item \textbf{Dokumentacja poprzez kod:} Interfejsy i typy jasno definiują strukturę danych przesyłanych między ekranami aplikacji.
\item \textbf{Refaktoryzacja:} Dzięki silnemu typowaniu, wprowadzanie zmian w dużych strukturach danych (np. modelu użytkownika) jest bezpieczniejsze i szybsze.
\end{itemize}

\subsection{Platforma Expo i usługi EAS}

Projekt oparto na środowisku \textbf{Expo} – ekosystemie narzędzi ułatwiających proces wytwarzania aplikacji w React Native. Expo dostarcza ujednolicone API do funkcji systemowych (aparat, galeria zdjęć, system plików). Wykorzystano również \textbf{EAS (Expo Application Services)}, co umożliwiło generowanie plików instalacyjnych (\textit{.apk} dla systemu Android) w chmurze, bez konieczności posiadania wydajnej stacji roboczej skonfigurowanej pod natywny proces kompilacji.

\subsection{Stylizacja interfejsu — NativeWind}

Do budowy warstwy wizualnej wykorzystano bibliotekę \textbf{NativeWind}, która implementuje system klas narzędziowych \textbf{Tailwind CSS} w środowisku mobilnym. Podejście to (\textit{utility-first}) pozwala na szybkie stylowanie komponentów bezpośrednio w kodzie JSX, co drastycznie skraca czas tworzenia responsywnych widoków i zapewnia spójność wizualną całego systemu bez konieczności pisania zewnętrznych arkuszy stylów.

\subsection{Nawigacja — React Navigation}

Zarządzanie przepływem ekranów zrealizowano przy pomocy biblioteki \textbf{React Navigation}. W systemie zaimplementowano dwa rodzaje nawigatorów:
\begin{itemize}
\item \textbf{Stack Navigator:} Wykorzystywany do liniowych procesów, takich jak logowanie, rejestracja czy przejście do szczegółów konkretnego rachunku.
\item \textbf{Bottom Tab Navigator:} Stanowi główne menu aplikacji, umożliwiając szybkie przełączanie się między modułami: Dom, Zadania, Finanse i Profil.
\end{itemize}

\subsection{Zarządzanie stanem globalnym — Zustand}

Do przechowywania danych dostępnych w całej aplikacji (np. informacja o zalogowanym użytkowniku, identyfikator wybranego domu) wybrano bibliotekę \textbf{Zustand}. Jest to rozwiązanie minimalistyczne, oparte na \textit{hookach}, które oferuje wysoką wydajność przy znacznie mniejszej złożoności niż tradycyjne biblioteki typu Redux. Pozwala to na uniknięcie problemu \textit{prop drilling} (przekazywania danych przez wiele poziomów komponentów).

\subsection{Komunikacja sieciowa — Axios}

Interakcja z backendem realizowana jest za pomocą klienta \textbf{Axios}. Biblioteka ta wspiera mechanizm \textit{interceptors}, który w projekcie został wykorzystany do automatycznego wstrzykiwania tokena JWT do nagłówków każdego żądania HTTP oraz do globalnej obsługi błędów sieciowych (np. automatyczne wylogowanie w przypadku wygaśnięcia sesji).

\subsection{Bezpieczeństwo — Expo Secure Store}

Ze względu na wrażliwy charakter danych dostępowych, w projekcie zrezygnowano z domyślnego \textit{AsyncStorage} na rzecz \textbf{Expo Secure Store}. Biblioteka ta zapewnia szyfrowane przechowywanie par klucz-wartość na poziomie systemowym (Keychain na iOS, SharedPreferences z szyfrowaniem na Androidzie), co gwarantuje, że tokeny autoryzacyjne użytkownika są chronione przed niepowołanym dostępem innych aplikacji.

\section{Integracja backend-frontend (REST API)}

Komunikacja między aplikacją mobilną a serwerem odbywa się zgodnie z zasadami architektury \textbf{REST} (Representational State Transfer).
Główne założenia REST spełnione w projekcie:
\begin{itemize}
    \item \textbf{Bezstanowość (Stateless):} Każde żądanie HTTP zawiera wszystkie informacje niezbędne do jego obsłużenia (np. token autoryzacyjny). Serwer nie przechowuje stanu sesji użytkownika między zapytaniami.
    \item \textbf{Jednolity interfejs:} Zasoby są identyfikowane przez adresy URL (np. \texttt{/api/homes}), a operacje na nich wykonywane są za pomocą standardowych metod HTTP (GET, POST, PUT, DELETE).
    \item \textbf{Format danych:} Wymiana danych odbywa się w lekkim formacie JSON, który jest natywnie wspierany zarówno przez JavaScript (Frontend), jak i biblioteki Go (Backend).
\end{itemize}

Wybrano REST zamiast GraphQL ze względu na prostotę implementacji, łatwość cachowania zapytań HTTP oraz jasny podział odpowiedzialności.

\section{Testy jednostkowe i integracyjne w Go}

Jakość kodu serwerowego jest weryfikowana za pomocą testów. Wykorzystano bibliotekę \textbf{Testify}, która rozszerza standardowy pakiet \texttt{testing} o asercje (np. \texttt{assert.Equal}, \texttt{assert.Nil}) oraz narzędzia do mockowania zależności. Testy obejmują kluczowe elementy logiki biznesowej oraz handlery HTTP.

\section{Podsumowanie}

Dobór technologii w projekcie ,,Household Manager'' podyktowany był chęcią stworzenia nowoczesnego, skalowalnego i łatwego w utrzymaniu systemu. Połączenie wydajnego backendu w Go z elastycznym frontendem w React Native stanowi solidną podstawę dla aplikacji użytkowej.
 % chapter3.tex zawiera treść rozdziału 3

%----------------------------------------------------------------------------------------
%	ROZDZIAŁ 4
%----------------------------------------------------------------------------------------
\chapter{Implementacja systemu}
\label{rozdzial4}

W tym rozdziale przedstawiono szczegóły implementacyjne systemu ,,Household Manager'', w tym strukturę kodu, konfigurację routingu, model danych, system autoryzacji oraz budowę aplikacji klienckiej. Zaprezentowano również kluczowe fragmenty kodu źródłowego oraz zrzuty ekranu gotowej aplikacji.

\section{Struktura projektu backend}

Aplikacja serwerowa została zaprojektowana zgodnie z zasadami architektury \textit{Clean Architecture} (Czysta Architektura), co zapewnia separację warstw, niezależność od frameworków oraz łatwość testowania. Kod podzielony jest na wyraźne warstwy odpowiedzialności:

\begin{enumerate}
    \item \textbf{Warstwa Transportowa (Handlers/Controllers):} Odpowiada za obsługę protokołu HTTP. Odbiera żądania, waliduje dane wejściowe (DTO), wywołuje odpowiednie serwisy i zwraca odpowiedź w formacie JSON. Nie zawiera logiki biznesowej.
    \item \textbf{Warstwa Usług (Services):} Serce aplikacji. Tutaj znajduje się cała logika biznesowa (np. obliczanie podziału rachunku, generowanie kodu zaproszenia). Serwisy operują na modelach domenowych i korzystają z repozytoriów.
    \item \textbf{Warstwa Dostępu do Danych (Repositories):} Odpowiada za bezpośrednią komunikację z bazą danych. Ukrywa szczegóły implementacyjne zapytań SQL/GORM przed resztą aplikacji.
    \item \textbf{Modele (Models):} Struktury danych reprezentujące encje biznesowe (np. User, Bill).
\end{enumerate}

Struktura katalogów projektu prezentuje się następująco:

\begin{itemize}
    \item \texttt{cmd/server/main.go}: Punkt wejścia aplikacji, inicjalizacja konfiguracji, połączenia z DB i start serwera.
    \item \texttt{internal/config}: Obsługa zmiennych środowiskowych (ładowanie z pliku .env).
    \item \texttt{internal/http/handlers}: Kontrolery HTTP.
    \item \texttt{internal/http/middleware}: Oprogramowanie pośredniczące (logowanie, autoryzacja JWT, obsługa CORS).
    \item \texttt{internal/models}: Definicje struktur danych i tagów GORM.
    \item \texttt{internal/repository}: Warstwa DAO (Data Access Object).
    \item \texttt{internal/services}: Logika biznesowa.
    \item \texttt{internal/router}: Konfiguracja routera Chi i mapowanie ścieżek.
\end{itemize}

\section{Endpoints API i routing w Chi}

Do obsługi routingu wykorzystano bibliotekę \texttt{go-chi}. Definiuje ona drzewiastą strukturę ścieżek, co ułatwia zarządzanie dużą liczbą endpointów.

\begin{lstlisting}[language=go, caption=Konfiguracja routingu dla zasobów Home, label=code:router]
r.Route("/homes", func(r chi.Router) {
    r.Post("/create", homeHandler.Create) // Tworzenie domu
    r.Post("/join", homeHandler.Join)     // Dolaczanie do domu
    r.Get("/my", homeHandler.GetUserHome) // Pobieranie domu uzytkownika

    r.Route("/{home_id}", func(r chi.Router) {
        // Middleware sprawdzajacy czy uzytkownik jest czlonkiem domu
        r.With(middleware.RequireMember(homeRepo)).Get("/", homeHandler.GetByID)
        
        // Zagniezdzone zasoby: zadania, rachunki, zakupy
        r.Route("/tasks", func(r chi.Router) {
            r.Get("/", taskHandler.GetTasksByHomeID)
            r.Post("/", taskHandler.Create)
        })
        r.Route("/bills", func(r chi.Router) {
             r.Get("/", billHandler.GetByHomeID)
        })
    })
})
\end{lstlisting}

\subsection{Dokumentacja API (Swagger)}

Wszystkie endpointy zostały udokumentowane przy użyciu standardu OpenAPI (Swagger). Poniższy zrzut ekranu przedstawia interaktywną dokumentację wygenerowaną automatycznie.

\begin{figure}[h]
    \centering
    
\includegraphics[width=0.95\textwidth]{placeholder.jpg} % Placeholder for Swagger UI
    \caption{Dokumentacja API w Swagger UI}
    \label{fig:swagger}
\end{figure}

\section{Modele danych i migracje}

Baza danych jest zarządzana przez bibliotekę GORM. Struktury Go są mapowane na tabele w bazie danych. Wykorzystano tagi strukturalne (\texttt{`gorm:"..."`}) do definiowania ograniczeń (constraints), indeksów oraz relacji.

Przykładowa, rozbudowana struktura modelu rachunku (\texttt{Bill}):

\begin{lstlisting}[language=go, caption=Model rachunku (Bill) w GORM, label=code:bill_model]
type Bill struct {
    ID          uuid.UUID `gorm:"type:uuid;default:gen_random_uuid();primaryKey"`
    Title       string    `gorm:"not null"`
    Amount      float64   `gorm:"not null"`
    PayerID     uuid.UUID `gorm:"type:uuid;not null"` // Klucz obcy do User
    Payer       User      `gorm:"foreignKey:PayerID"` // Relacja
    HomeID      uuid.UUID `gorm:"type:uuid;not null"` // Klucz obcy do Home
    CategoryID  *uuid.UUID `gorm:"type:uuid"`
    Category    *BillCategory `gorm:"foreignKey:CategoryID"`
    DueDate     time.Time
    IsPaid      bool      `gorm:"default:false"`
    ReceiptURL  string    // URL do zdjecia w S3
    
    // Pola audytowe
    CreatedAt   time.Time
    UpdatedAt   time.Time
    DeletedAt   gorm.DeletedAt `gorm:"index"` // Obsluga Soft Delete
}
\end{lstlisting}

\section{System autoryzacji i JWT}

Bezpieczeństwo API opiera się na tokenach JWT (JSON Web Tokens). Proces autoryzacji przebiega następująco:
\begin{enumerate}
    \item Użytkownik loguje się podając email i hasło.
    \item Serwer weryfikuje hasło (porównuje hash bcrypt).
    \item W przypadku sukcesu, generowany jest token JWT zawierający \texttt{user\_id} oraz datę wygaśnięcia (np. 72h).
    \item Token jest podpisywany tajnym kluczem (\texttt{JWT\_SECRET}) przechowywanym tylko na serwerze.
    \item Klient dołącza token do nagłówka \texttt{Authorization: Bearer <token>} każdego żądania.
\end{enumerate}

\section{Obsługa błędów i walidacja danych}

Wszystkie błędy zwracane przez API mają ujednoliconą strukturę JSON, co ułatwia ich obsługę po stronie klienta. Walidacja danych wejściowych odbywa się przy użyciu biblioteki \texttt{go-playground/validator}.

\begin{lstlisting}[language=json, caption=Struktura błędu API, label=json:error]
{
    "error": "Validation Failed",
    "details": {
        "Email": "field is required",
        "Password": "min length is 8"
    }
}
\end{lstlisting}

\section{Integracja z OCR dla rachunków}

Moduł OCR (\textit{Optical Character Recognition}) umożliwia automatyczne zaczytywanie danych z rachunków. Logika parsowania tekstu z paragonu opiera się na wyrażeniach regularnych (Regex), które poszukują wzorców dat i kwot.

Fragment serwisu odpowiedzialnego za analizę tekstu z paragonu:

\begin{lstlisting}[language=go, caption=Logika parsowania tekstu z paragonu, label=code:ocr_logic]
func (s *OCRService) ParseReceiptText(text string) (float64, string) {
    // Regex dla kwot (np. 123.45, 123,45 PLN)
    amountRegex := regexp.MustCompile(`(\d+[.,]\d{2})\s*(?:PLN|zl|EUR)?`)
    
    lines := strings.Split(text, "\n")
    var maxAmount float64
    var vendor string

    for i, line := range lines {
        if i == 0 {
            vendor = line // Zalozenie: pierwsza linia to nazwa sklepu
        }
        matches := amountRegex.FindStringSubmatch(line)
        if len(matches) > 1 {
            // Konwersja przecinka na kropke i parsowanie float
            val, err := strconv.ParseFloat(strings.Replace(matches[1], ",", ".", 1), 64)
            if err == nil && val > maxAmount {
                maxAmount = val // Szukamy najwiekszej kwoty (Suma)
            }
        }
    }
    return maxAmount, vendor
}
\end{lstlisting}

\section{Struktura projektu frontend}

Aplikacja kliencka w React Native wykorzystuje \textbf{Expo Router}, który opiera routing na strukturze plików (File-based routing), podobnie jak framework Next.js.
Katalog \texttt{app} zawiera ekrany aplikacji:

\begin{itemize}
    \item \texttt{app/index.tsx}: Ekran startowy (Splash screen).
    \item \texttt{app/(auth)/login.tsx}: Ekran logowania (grupa auth, niedostępna po zalogowaniu).
    \item \texttt{app/(tabs)/\_layout.tsx}: Konfiguracja dolnego paska nawigacji (Tab Bar).
    \item \texttt{app/(tabs)/home.tsx}: Główny pulpit użytkownika.
    \item \texttt{app/[id].tsx}: Dynamiczna ścieżka (np. szczegóły zadania).
\end{itemize}

\section{Interfejs użytkownika (UI) - Prezentacja aplikacji}

Interfejs użytkownika został zaprojektowany z naciskiem na przejrzystość i łatwość obsługi (User Experience). Poniżej przedstawiono kluczowe ekrany aplikacji.

\subsection{Ekran logowania i rejestracji}
Umożliwia dostęp do aplikacji. Użytkownik może zalogować się tradycyjnie lub przez Google.

\begin{figure}[h]
    \centering
    
\includegraphics[width=0.45\textwidth]{placeholder.jpg} % Placeholder: Login Screen
    \caption{Ekran logowania}
    \label{fig:login_screen}
\end{figure}

\subsection{Pulpit główny (Dashboard)}
Centrum dowodzenia domem. Wyświetla podsumowanie: najbliższe zadania, stan konta oraz skróty do najważniejszych funkcji.

\begin{figure}[h]
    \centering
    
\includegraphics[width=0.45\textwidth]{placeholder.jpg} % Placeholder: Dashboard
    \caption{Pulpit główny aplikacji}
    \label{fig:dashboard_screen}
\end{figure}

\subsection{Zarządzanie zadaniami}
Lista zadań z możliwością filtrowania i oznaczania jako wykonane.

\begin{figure}[h]
    \centering
    
\includegraphics[width=0.45\textwidth]{placeholder.jpg} % Placeholder: Tasks
    \caption{Lista zadań domowych}
    \label{fig:tasks_screen}
\end{figure}

\subsection{Skanowanie rachunków (OCR)}
Moduł pozwalający na zrobienie zdjęcia i weryfikację odczytanych danych.

\begin{figure}[h]
    \centering
    
\includegraphics[width=0.45\textwidth]{placeholder.jpg} % Placeholder: OCR/Bill
    \caption{Ekran dodawania rachunku ze skanowaniem OCR}
    \label{fig:ocr_screen}
\end{figure}

\section{Wybrane fragmenty kodu frontendowego}

Poniżej przedstawiono niestandardowy React Hook odpowiedzialny za pobieranie danych z API z obsługą stanu ładowania i błędów.

\begin{lstlisting}[language=javascript, caption=Custom Hook useFetch, label=code:usefetch]
export function useFetch<T>(url: string) {
  const [data, setData] = useState<T | null>(null);
  const [loading, setLoading] = useState(true);
  const [error, setError] = useState(null);

  useEffect(() => {
    const fetchData = async () => {
      try {
        const response = await api.get(url);
        setData(response.data);
      } catch (err) {
        setError(err);
      } finally {
        setLoading(false);
      }
    };
    fetchData();
  }, [url]);

  return { data, loading, error, refetch: fetchData };
}
\end{lstlisting}

 % chapter4.tex zawiera treść rozdziału 4

%----------------------------------------------------------------------------------------
%	ROZDZIAŁ 5
%----------------------------------------------------------------------------------------
\chapter{Badania i testy systemu}
\label{rozdzial5}

Celem tego rozdziału jest weryfikacja poprawności działania wytworzonego oprogramowania oraz ocena stopnia realizacji założonych wymagań. Przeprowadzono szereg testów funkcjonalnych oraz wydajnościowych.

\section{Scenariusze testowe}

W celu sprawdzenia kluczowych funkcjonalności systemu opracowano następujące scenariusze testowe, odzwierciedlające typowe ścieżki użytkownika:

\subsection{Scenariusz 1: Konfiguracja nowego gospodarstwa domowego}
\begin{enumerate}
    \item Użytkownik A rejestruje nowe konto.
    \item Użytkownik A tworzy nowy Dom o nazwie ,,Mieszkanie Studenckie''.
    \item Użytkownik A generuje kod zaproszenia i przekazuje go Użytkownikowi B.
    \item Użytkownik B rejestruje się i dołącza do domu używając otrzymanego kodu.
    \item \textbf{Oczekiwany rezultat:} Obaj użytkownicy widzą ten sam pulpit domu i mają dostęp do wspólnych zasobów.
\end{enumerate}

\subsection{Scenariusz 2: Dodawanie i rozliczanie rachunku z OCR}
\begin{enumerate}
    \item Użytkownik wybiera opcję dodania rachunku i robi zdjęcie paragonu.
    \item System przetwarza obraz i proponuje kwotę oraz nazwę sklepu.
    \item Użytkownik zatwierdza dane i wybiera podział kosztów ,,po równo''.
    \item Rachunek pojawia się na liście jako ,,nieopłacony''.
    \item \textbf{Oczekiwany rezultat:} Kwota zostaje poprawnie odczytana, a dług zostaje przypisany do pozostałych domowników.
\end{enumerate}

\section{Testy funkcjonalne backendu}

Testy jednostkowe i integracyjne warstwy serwerowej zostały zaimplementowane przy użyciu pakietu \texttt{testing} oraz biblioteki \texttt{testify}.

Przykładowy wynik uruchomienia testów dla serwisu autoryzacji:
\begin{lstlisting}[language=bash, caption=Wynik testów jednostkowych, label=test:output]
$ go test ./internal/services/... -v
=== RUN   TestRegisterUser
--- PASS: TestRegisterUser (0.00s)
=== RUN   TestLoginUser_Success
--- PASS: TestLoginUser_Success (0.00s)
=== RUN   TestLoginUser_InvalidPassword
--- PASS: TestLoginUser_InvalidPassword (0.00s)
PASS
ok      github.com/Dragodui/diploma-server/internal/services    0.123s
\end{lstlisting}

Wszystkie kluczowe ścieżki krytyczne (rejestracja, tworzenie domu, dodawanie zadań) zostały pokryte testami automatycznymi, co zapewnia stabilność API.

\section{Testy interfejsu mobilnego}

Testy aplikacji mobilnej przeprowadzono manualnie na dwóch urządzeniach fizycznych:
\begin{itemize}
    \item Google Pixel 6 (Android 14)
    \item iPhone 13 (iOS 17)
\end{itemize}
oraz na emulatorach.

Weryfikowano:
\begin{itemize}
    \item \textbf{Responsywność:} Czy elementy UI skalują się poprawnie na różnych rozdzielczościach.
    \item \textbf{Nawigację:} Czy przechodzenie między ekranami jest płynne i intuicyjne.
    \item \textbf{Obsługę błędów:} Czy aplikacja wyświetla czytelne komunikaty w przypadku braku internetu lub błędu serwera.
\end{itemize}

Stwierdzono, że aplikacja działa stabilnie na obu platformach. Wykryte drobne błędy wizualne (np. przesunięcia tekstu) zostały na bieżąco poprawione.

\section{Wyniki i wnioski z testów}

Przeprowadzone badania wykazały, że system spełnia założone wymagania funkcjonalne.

\begin{itemize}
    \item \textbf{Wydajność:} Czas odpowiedzi API dla standardowych operacji wynosi średnio poniżej 100ms (w środowisku lokalnym).
    \item \textbf{Skuteczność OCR:} Moduł rozpoznawania tekstu poprawnie odczytuje kwoty z paragonów w około 80% przypadków przy dobrym oświetleniu. W przypadku słabej jakości zdjęcia, użytkownik ma możliwość ręcznej edycji danych.
    \item \textbf{Użyteczność:} Testy z udziałem użytkowników potwierdziły, że interfejs jest przejrzysty, a proces zakładania domu i zapraszania domowników nie sprawia trudności.
\end{itemize}

System jest gotowy do wdrożenia w wersji produkcyjnej (MVP).
 % chapter5.tex zawiera treść rozdziału 5

%----------------------------------------------------------------------------------------
%	ROZDZIAŁy kolejne należy dodać analogicznie do 1 i 2 
% utworzyć pliki i je załączyć (include)
%----------------------------------------------------------------------------------------

%----------------------------------------------------------------------------------------
%	ZAKOŃCZENIE PRACY DYPLOMOWEJ (JK - design and implementation)
%----------------------------------------------------------------------------------------
\addcontentsline{toc}{chapter}{Podsumowanie}
\chapter{Podsumowanie i wnioski}
\label{rozdzial6}

Celem niniejszej pracy inżynierskiej było zaprojektowanie i zaimplementowanie systemu wspomagającego zarządzanie gospodarstwem domowym. Osiągnięto to poprzez stworzenie aplikacji mobilnej ,,Household Manager'' oraz dedykowanego serwera API.

\section{Realizacja celów pracy}

Wszystkie główne założenia projektowe zostały zrealizowane. Powstała aplikacja umożliwia:
\begin{itemize}
    \item Tworzenie grup użytkowników (domów) i zarządzanie ich członkami.
    \item Ewidencjonowanie i przydzielanie zadań domowych.
    \item Transparentne zarządzanie finansami i rozliczanie wspólnych wydatków.
    \item Automatyzację wprowadzania danych z rachunków dzięki technologii OCR.
    \item Wspólne planowanie zakupów i demokratyczne podejmowanie decyzji poprzez ankiety.
\end{itemize}

Wykorzystanie nowoczesnych technologii, takich jak język Go na backendzie oraz React Native na frontendzie, pozwoliło na stworzenie wydajnego i skalowalnego rozwiązania, działającego na dwóch najpopularniejszych platformach mobilnych (Android i iOS).

\section{Wnioski z realizacji projektu}

Podczas realizacji projektu wyciągnięto następujące wnioski:
\begin{itemize}
    \item Architektura mikroserwisowa (lub modularny monolit w kontenerach Docker) znacząco ułatwia wdrażanie i skalowanie aplikacji.
    \item Biblioteka GORM znacznie przyspiesza pracę z bazą danych, jednak w przypadku bardzo skomplikowanych zapytań SQL konieczna jest znajomość natywnego języka zapytań.
    \item Technologia Expo w ekosystemie React Native jest dojrzałym rozwiązaniem, które pozwala na szybkie prototypowanie i budowanie produkcyjnych aplikacji bez konieczności głębokiej ingerencji w kod natywny (Java/Swift).
    \item Integracja z systemami OCR jest wyzwaniem ze względu na różnorodność formatów paragonów i jakość zdjęć. Wymaga to implementacji mechanizmów walidacji i korekty danych przez użytkownika.
\end{itemize}

\section{Kierunki dalszego rozwoju}

System posiada potencjał do dalszej rozbudowy. Możliwe kierunki rozwoju to:
\begin{enumerate}
    \item \textbf{Wersja przeglądarkowa (Web):} Stworzenie panelu administracyjnego dostępnego przez przeglądarkę internetową.
    \item \textbf{Integracja z bankowością (Open Banking):} Automatyczne pobieranie transakcji z kont bankowych użytkowników w celu łatwiejszego rozliczania.
    \item \textbf{Rozbudowa modułu AI:} Wykorzystanie bardziej zaawansowanych modeli uczenia maszynowego do lepszego kategoryzowania wydatków i inteligentnego przydzielania zadań (np. na podstawie preferencji i historii).
    \item \textbf{Kalendarz:} Wizualizacja zadań i płatności na wspólnym kalendarzu z możliwością synchronizacji z Google Calendar.
\end{enumerate}

Podsumowując, aplikacja ,,Household Manager'' stanowi funkcjonalne i nowoczesne narzędzie, które realnie może wpłynąć na poprawę jakości życia osób dzielących gospodarstwo domowe, redukując liczbę konfliktów i usprawniając codzienne obowiązki.
 % conclusions.tex zawiera treść zakończenia/podsumowania pracy dyplomowej


%----------------------------------------------------------------------------------------
%	SPIS LITERATURY (JK - design and implementation)
%----------------------------------------------------------------------------------------
\pagestyle{plain}
\begin{thebibliography}{99}
\addcontentsline{toc}{chapter}{Spis literatury}

\bibitem{donovan2015} Donovan A. A. A., Kernighan B. W., \textit{The Go Programming Language}, Addison-Wesley, New York 2015.
\bibitem{cox2016} Cox M., Campoy F., \textit{Go in Practice}, Manning Publications, Shelter Island 2016.
\bibitem{freeman2020} Freeman E., Robson E., \textit{Head First Design Patterns}, O'Reilly Media, Sebastopol 2020.
\bibitem{evans2004} Evans E., \textit{Domain-Driven Design: Tackling Complexity in the Heart of Software}, Addison-Wesley, Boston 2004.
\bibitem{fowler2002} Fowler M., \textit{Patterns of Enterprise Application Architecture}, Addison-Wesley, Boston 2002.
\bibitem{martin2018} Martin R. C., \textit{Clean Architecture: A Craftsman's Guide to Software Structure and Design}, Prentice Hall, Boston 2018.
\bibitem{pilone2005} Pilone D., \textit{UML 2.0 in a Nutshell}, O'Reilly Media, Sebastopol 2005.
\bibitem{date2003} Date C. J., \textit{An Introduction to Database Systems}, Addison-Wesley, Boston 2003.
\bibitem{burrows2020} Burrows M., \textit{The Go Memory Model}, Google Research, 2020.
\bibitem{gorm2023} GORM Team, \textit{Design and Best Practices for ORM in Go}, Medium, 2023.
\bibitem{rn2024} React Native Team, \textit{State Management in React Native Apps}, Expo Blog, 2024.
\bibitem{postgres2023} PostgreSQL Global Development Group, \textit{Indexing and Query Optimization in PostgreSQL}, PostgreSQL Tech Notes, 2023.
\bibitem{godocs} Go Team, \textit{Go Documentation}, \url{https://golang.org/doc/}, 2025.
\bibitem{gormdocs} GORM, \textit{GORM Official Documentation}, \url{https://gorm.io/docs/}, 2025.
\bibitem{chidocs} Go-Chi, \textit{Chi Router Documentation}, \url{https://github.com/go-chi/chi}, 2025.
\bibitem{rndocs} Meta, \textit{React Native Documentation}, \url{https://reactnative.dev/docs/}, 2025.
\bibitem{expodocs} Expo, \textit{Expo Documentation}, \url{https://docs.expo.dev/}, 2025.
\bibitem{zustanddocs} Poimandres, \textit{Zustand Documentation}, \url{https://docs.pmnd.rs/zustand}, 2025.
\bibitem{axiosdocs} Axios, \textit{Axios Documentation}, \url{https://axios-http.com/docs/intro}, 2025.
\bibitem{pgdocs} PostgreSQL, \textit{PostgreSQL Official Documentation}, \url{https://www.postgresql.org/docs/}, 2025.
\bibitem{dockerdocs} Docker Inc., \textit{Docker Documentation}, \url{https://docs.docker.com/}, 2025.
\bibitem{gitdocs} Git, \textit{Git Documentation}, \url{https://git-scm.com/doc}, 2025.
\bibitem{owasp2024} OWASP Foundation, \textit{OWASP Top 10: Security Risks for Web Applications}, \url{https://owasp.org/www-project-top-ten/}, 2024.
\bibitem{tesseract} Tesseract OCR, \textit{Tesseract OCR Documentation}, \url{https://tesseract-ocr.github.io/}, 2024.
\bibitem{googledev} Google Developers, \textit{Push Notifications for Mobile Apps}, \url{https://firebase.google.com/docs/cloud-messaging}, 2024.

\end{thebibliography}

%----------------------------------------------------------------------------------------
% Koniec dokumentu
\end{document}
